% Generated draft: requires mathematical and source review.
\documentclass[11pt]{article}
\usepackage[margin=25mm]{geometry}
\usepackage{fontspec}
\setmainfont{Latin Modern Roman}
\usepackage{amsmath,amssymb,mathtools,unicode-math}
\setmathfont{Latin Modern Math}
\usepackage{enumitem,needspace,xurl,hyperref}
\setlength{\parindent}{0pt}
\setlength{\parskip}{5pt}
\setlength{\emergencystretch}{3em}
\newcommand{\N}{\mathbb{N}}
\newcommand{\Z}{\mathbb{Z}}
\newcommand{\Q}{\mathbb{Q}}
\newcommand{\R}{\mathbb{R}}
\newcommand{\C}{\mathbb{C}}
\newcommand{\F}{\mathbb{F}}
\newcommand{\A}{\mathbb{A}}
\newcommand{\PP}{\mathbb{P}}
\newcommand{\E}{\mathbb{E}}
\newcommand{\eps}{\varepsilon}
\newcommand{\dd}{\,\mathrm{d}}
\begin{document}
\section*{UnsolvedMath ID: 800475}
\textbf{Source problem code:} catalogue entry
\subsection{Definitions and mathematical statement}
A polygonal tiling billiard is the directional dynamical system associated with a polygonal tiling \(\mathcal T\) of \(\R^2\). A particle follows a straight segment inside a tile and, upon crossing an edge into an adjacent tile, changes direction according to Snell's law with refraction coefficient \(-1\). With speed normalized to one, the usual negative-refraction convention at a regular edge crossing is
\[
 v^+=2\langle v^-,n\rangle n-v^-,
\]
where \(n\) is a unit normal pointing into the next tile: the normal component is preserved and the tangential component is reversed. The supplied record does not specify continuation through vertices or tangential contacts, so these cases remain outside this description.

The source calls behavior \emph{integrable} when trajectories are typically either periodic or escape linearly to infinity. It does not specify the meaning of ``typically,'' the relevant parameter and trajectory measures, or a precise definition of linear or nonlinear escape. Its family-by-family account is as follows; these are supplied descriptions and literature attributions, not independently verified assertions.
\begin{enumerate}
\item The \emph{trihexagonal tiling} is the periodic tiling by regular hexagons and triangles. The source attributes to Diana Davis and Pat Hooper [DH18] the existence of the described orbit behavior: density in an unbounded periodic open set. The quantifier over trajectories is not made precise.
\item A \emph{triangle tiling} is a periodic tiling generated from one triangular tile by central symmetries. The source describes its behavior as integrable, with a measure-zero exceptional set of triangle parameters, parametrized by the Rauzy gasket, admitting nonlinearly escaping trajectories. It states that these exceptional trajectories necessarily pass through the circumcenters of every tile they cross. These claims are attributed to [BDFI18, TAF, H\_Paro], with [highlights\_O] as further discussion. The source describes the Rauzy gasket as an attractor of a multidimensional continued-fraction algorithm and refers to [AR91, AHS, DDe, AS]; it supplies neither the algorithm nor the parametrization or measure. It also recommends Ofir David's animation [Ofir], based on [TAF].
\item A \emph{cyclic quadrilateral tiling} is a periodic tiling generated from one cyclic quadrilateral by central symmetries. The source describes its behavior as integrable [H\_Paro, TAF], while also reporting exceptional tilings admitting nonlinearly escaping trajectories. It attributes their classification and parametrization by Novikov's gasket to [Dynnikovetal]. Thus the integrable description must retain its ``typically'' qualification; the source supplies no precise genericity convention or definition of this gasket.
\item A \emph{double triangle tiling} is a periodic tiling by a pair of triangles \(\Delta,\Delta'\) sharing an edge \(a\), whose opposite angles satisfy \(\alpha=\alpha'\). The source describes these systems as generalizations of triangle tiling billiards, attributes partial dynamical results to Magali Jay's thesis [Jay\_thesis], and suggests, without asserting a theorem, that their behavior is probably integrable. The complete periodic assembly and parameter space are not specified.
\item The \emph{refractive wind-tree model} uses rectangles of size \(a\times b\) whose centers belong to a lattice \(\Lambda\), together with the complement of their union, described by the source as an ``unbounded polygon.'' The source reports that generic trajectories are trapped in bands [Jay\_thesis]. Rectangle orientation, nonoverlap conditions, the exact placement rule, and the meaning of genericity and of a trapping band are not specified.
\item A \emph{generalized tiling billiard} starts with a fixed polygon \(P\), whether or not it tiles the plane. After a trajectory crosses a side, the next copy of \(P\) is obtained by central symmetry about that side's midpoint; the trajectory continues with local refraction coefficient \(-1\). This describes successive polygon copies along a trajectory, without asserting that they form a global tiling. For cyclic pentagons, the source reports nonlinear escape in open sets, with a principal direction and deviations from that direction [Jay\_pentagons, Jay\_thesis]. It states that generic behavior is unknown. The ambient space of the reported open sets and the precise meanings of principal direction and deviations are unspecified.
\item A \emph{brick tiling} is described as a periodic tiling formed from two rows of \(1\times1\) squares shifted relative to one another by \(\theta\in(0,1)\). The source attributes periodicity of all orbits for \(\theta\in\Q\) to [thesis\_park] and states that behavior for \(\theta\notin\Q\) has not yet been studied. The extension of the two-row description to the entire plane and conventions for singular orbits are not supplied.
\end{enumerate}

The source calls a polygonal tiling \emph{locally foldable} if its tiles admit a two-coloring and the alternating sum of incident tile angles at each vertex is zero. Under the usual interpretation that adjacent tiles have opposite colors, if the incident angles in cyclic order are \(\beta_1,\ldots,\beta_{2m}\), this defect condition is
\[
\sum_{j=1}^{2m}(-1)^{j-1}\beta_j=0.
\]
The source identifies families 2, 3, and 4 as locally foldable. It attributes to [TAF] the statements that all trajectories in such tilings are either periodic or escaping and that any trajectory passes through each tile at most once. Taken literally over an entire trajectory, the latter statement conflicts with periodic return; the intended restriction or convention is unresolved. The source relates locally foldable tiling billiards to Novikov's problem in topology, referring to [Dynnikov\_main] for an introduction and [PaRoH, Dynnikovetal] for the connection. It further states that all the cited tiling-billiard work reduces the dynamics to linear flows on families of translation or half-translation surfaces, with constructions specific to each family.

\emph{Unresolved problem statement.} The supplied record is an overview bearing the title ``Olga Paris-Romaskevich --- Problem 1'' but contains no explicit single question or quantified conjecture. Its stated unresolved directions include possible integrability of double triangle tiling billiards, generic dynamics for cyclic pentagonal generalized tiling billiards, and dynamics of brick tiling billiards for irrational \(\theta\). These directions are retained as source-reported research questions; choosing one as the intended target, defining genericity, or imposing missing family hypotheses would require information absent from the record.
\subsection{Short English statement}
This record surveys polygonal tiling billiards with refraction coefficient \(-1\), attributes known and exceptional behavior to supplied references, and identifies several unresolved directions. It supplies no single precise question. The proposed integrability of double triangle systems, generic cyclic pentagon dynamics, and irrational brick dynamics therefore remain separate, incompletely specified targets. The registry labels the record open; this draft requires human mathematical and source review.
\subsection{Sources}
\begin{itemize}
\item Supplied local UnsolvedMath registry record: permanent ID \(800475\), exact source problem code \texttt{catalogue entry}, title ``Billiards and Quantitative Symplectic Geometry --- Olga Paris-Romaskevich --- Problem 1.'' The supplied registry status is \texttt{open}; it is not an independently verified literature assessment. The record lists Albach, Barraud et al. as proposers, year 2026, and Olga Paris-Romaskevich as contributor.
\item Supplied source collection: Albach, Barraud et al., \emph{Open Problems in Billiards and Quantitative Symplectic Geometry} (2026), Olga Paris-Romaskevich, Problem 1. Supplied URL: \url{https://arxiv.org/abs/2602.12896}. The record gives access date \texttt{2026-10-16} and extraction method \texttt{arxiv\_latex}. These are supplied provenance details; the source has not been independently visited or checked for this draft. The contributor qualifies the survey as a complete list only ``as far as I know,'' covering families studied since 2015.
\item Trihexagonal tilings: [DH18], attributed in the source to Diana Davis and Pat Hooper. Triangle tilings: [BDFI18], [TAF], [H\_Paro], and [highlights\_O]. Rauzy-gasket background: [AR91], [AHS], [DDe], and [AS]. Animation: [Ofir], attributed to Ofir David and described as based on [TAF].
\item Cyclic quadrilateral tilings: [H\_Paro], [TAF], and [Dynnikovetal]. Double triangle tilings and refractive wind-tree dynamics: [Jay\_thesis], identified as Magali Jay's thesis. Cyclic pentagonal generalized systems: [Jay\_pentagons] and [Jay\_thesis]. Brick tilings: [thesis\_park].
\item Local foldability: [TAF]. Novikov's problem and its connection to tiling billiards: [Dynnikov\_main], [PaRoH], and [Dynnikovetal]. These reference labels are reproduced from the supplied record; their full bibliographic entries and individual links were not supplied.
\end{itemize}
\end{document}
